\subsection*{Simplifications}

Every time the value numbering process generates a signature, it tries to simplify the expression. For example,

\begin{lstlisting}[language=croqtile,caption={A constant expression}]
   int a = 3 + 8;
\end{lstlisting}

Given the simple integer expression, the value numbering process could generate signatures and associated value number as illustrated:

\begin{lstlisting}[language=croqtile,caption={The simplified graph}]
   VN #0:  const_3
   VN #1:  const_8
   VN #2:  +:#0:#1
   Alias a -> #2
\end{lstlisting}

Here, without optimizations, two constant values \#0 and \#1, and a new value number \#2, which is an addition of \#0 and \#1 are generated. \#2 has the signature of "+:\#0:\#1". However, considering that addition of constants can be evaluated at compile time, simplification could be performed. The following output showcases the optimized value numbering process:

\begin{lstlisting}[language=croqtile,caption={The folded graph}]
   VN #0:  const_3
   VN #1:  const_8
   VN #2:  const_11
   Alias a -> #2
\end{lstlisting}

The signature of \#2 now represents a constant, which implies the algebraic simplification is applied successfully.
